\documentclass[10pt,a4paper]{amsart}
\usepackage[margin=19mm]{geometry}
\usepackage{amsmath,amssymb,mathtools}
\usepackage[scr=rsfs,cal=euler]{mathalfa}
\usepackage{xfrac}
\usepackage[unicode,hidelinks,bookmarks=false]{hyperref}
\newcommand{\ie}{i.e.\ }
\newcommand{\cf}{cf.\ }
\newcommand{\resp}{respectively\ }
\newcommand{\Subsection}{\subsection}
\providecommand{\nobreakdash}{\nobreak\hbox{-}}
\setlength{\emergencystretch}{2em}
\multlinegap0pt
\usepackage{bm}
\usepackage{bbm}
\usepackage{dsfont}
\usepackage{xspace}
\usepackage{cancel}
\usepackage{mathtools}
\usepackage{amsthm}
\makeatletter
\@ifundefined{Proposition}{\newtheorem{Proposition}{Proposition}}{}
\@ifundefined{Theorem}{\newtheorem{Theorem}{Theorem}}{}
\makeatother
\newcommand{\Ff}{{\mathbb F}}
\title[New bounds and constructions for large partial $m$-ovoids an]{New bounds and constructions for large partial $m$-ovoids and related structures}
\author{John Bamberg, Anurag Bishnoi, Ferdinand Ihringer, Ananthakrishnan Ravi}
\date{Source: arXiv:2406.03043v3 (CC BY 4.0)}
\begin{document}
\maketitle

\noindent\emph{Source and licence.} New bounds and constructions for large partial $m$-ovoids and related structures. Version arXiv:2406.03043v3 (CC BY 4.0), distributed under the Creative Commons Attribution 4.0 International licence, as recorded in the arXiv licence metadata for this exact version and shown on the abstract page https://arxiv.org/abs/2406.03043v3. Full rights notice: licenses/arxiv-2406.03043-CC-BY-4.0.txt.\par\smallskip

\noindent\textit{Editorial scope.} Counted unit: the Proposition whose source label is \texttt{low2rank} in the article (source0.tex lines 862--884, source label \texttt{low2rank}), delivered together with the article's own complete original proof of it (source0.tex lines 886--981, one \texttt{proof} environment of 96 lines). No mathematics is written, repaired or re-derived: statement and proof are the article's own text line for line. Required context retained uncounted: thm:BLS (lines 843--857). Every definition, display and label of those lines is retained unchanged. Omitted from this package and not redistributed: 1--842, 858--861, 885--885, 982--1234. Nothing inside a retained excerpt is omitted or altered: every retained body is delivered exactly as the article's lines are written, so no omission receipt of any kind accompanies this package. Citations: the retained excerpts carry no citation command at all, so no bibliography entry of the article is transcribed and no reference is redistributed. Reference adaptation: none was needed and none was made; every reference inside the delivered unit resolves inside it, so no unresolved reference or citation is produced. Printed slips kept literal: no printed word, symbol, hypothesis or number of the counted statement or of its proof was altered. Layout and class: the article's own class is \texttt{amsproc} with packages including amssymb, amsmath, amsthm, lmodern, tikzsymbols, amsaddr, framed, booktabs, float, xcolor, hyperref, enumerate, mathtools, url; the wrapper is the lane's pinned \texttt{amsart} editorial wrapper at 10pt on A4, which restates the theorem-like environments the retained text needs and adds the article's own macro definitions that the retained text needs (\texttt{\textbackslash Ff}), with extra wrapper packages (bm, bbm, dsfont, xspace, cancel, mathtools). The delivered numbering is the unit's own. Assets: no retained span contains a figure environment, a graphics-inclusion command, a PDF-page inclusion, a table or a TikZ picture; no image file is copied into this package, the package declares no asset dependency and no image file is redistributed with it. Licence: version arXiv:2406.03043v3 of this article is distributed under the Creative Commons Attribution 4.0 International licence, as recorded in the arXiv licence metadata for this exact version and shown on the abstract page https://arxiv.org/abs/2406.03043v3. A full rights notice is delivered at licenses/arxiv-2406.03043-CC-BY-4.0.txt. Characters: every retained excerpt is pure ASCII, so the delivered engine, XeLaTeX with its default Latin Modern text font, reports no missing character.
\begin{Theorem}\label{thm:BLS}
  Let $m > 1$ be a positive integer.
  Let $\Gamma$ be a graph on $n$ vertices and $r := \mathrm{rank}_\Ff(A_\Gamma+I)$
  for some field $\Ff$. Then, for every large enough $h$, there exists a graph $\Delta$
  such that:
  \begin{enumerate}
   \item the clique number of $\Delta$ is at most $m$,
   \item $\mathrm{rank}_\Ff(A_\Delta+I) \leq r^h$,
   \item the graph $\Delta$ has
   \[
    \frac{1}{m} \left( \frac{n}{\left( \sum_{t=1}^m t^{m+1} a_t \right)^{\frac{1}{m+1}} } \right)^h - 1
   \]
  vertices. Here $a_i$ denotes the number of cliques of size $i$ in $\Gamma$.
  \end{enumerate}
\end{Theorem}

\begin{Proposition}\label{low2rank}
For every fixed integer $m\geq 2$ and all sufficiently large $t$, there exists a
$K_{m+1}$-free graph with complementary $2$-rank exactly $t$ and at least
\[
t^{\beta_m-o(1)}
\]
vertices, where
\[
\beta_m :=
\log_{10}\left(
\frac{32}{\left(32+144\cdot 2^{m+1}\right)^{1/(m+1)}}
\right).
\]


Moreover, for every $\varepsilon>0$, there exists an integer $m$ such that, for
all sufficiently large $t$, there exists a $K_{m+1}$-free graph with
complementary $2$-rank exactly $t$ and at least
\[
t^{\frac{\log 16}{\log 10}-\varepsilon}
\]
vertices.
\end{Proposition}

\begin{proof}
Let $\Gamma$ be the Cayley graph constructed above for $n=5$. Then $\Gamma$ has
$32$ vertices and $144$ edges, and its spectrum is
\[
9^{(1)},\; 5^{(1)},\; 1^{(21)},\; (-3)^{(7)},\; (-7)^{(2)}.
\]
The multiplicity of the eigenvalue $1$ gives
\[
\operatorname{rank}_{\mathbb R}(A_\Gamma-I)=11.
\]
Since the binary rank of an integral matrix is at most its rank over
$\mathbb R$, we get
\[
\operatorname{rank}_{\mathbb F_2}(A_\Gamma+I)
=
\operatorname{rank}_{\mathbb F_2}(A_\Gamma-I)
\leq 11.
\]
In fact, a direct computation gives
\[
\operatorname{rank}_{\mathbb F_2}(A_\Gamma+I)=10.
\]

We now apply Theorem~\ref{thm:BLS} to this graph with
\[
a_1=32,\qquad a_2=144.
\]
For each fixed integer $m\geq 2$ and each positive integer $h$, the theorem
produces a $K_{m+1}$-free graph $\Delta_h$ with complementary $2$-rank at most
$10^h$ and at least
\[
\frac{1}{m+1}
\left(
\frac{32}{\left(32+144\cdot 2^{m+1}\right)^{1/(m+1)}}
\right)^h
\]
vertices.

Since
\[
10^h = T
\]
implies
\[
\left(
\frac{32}{\left(32+144\cdot 2^{m+1}\right)^{1/(m+1)}}
\right)^h
=
T^{\beta_m},
\]
we obtain, for infinitely many values $T=10^h$, a $K_{m+1}$-free graph with
complementary $2$-rank at most $T$ and at least
\[
T^{\beta_m-o(1)}
\]
vertices.

It remains only to pass from the sequence $T=10^h$ to every sufficiently large
$t$. Given $t$, choose $h$ such that
\[
10^h\leq t<10^{h+1}.
\]
Take the graph $\Delta_h$ and add isolated vertices until the complementary
$2$-rank becomes exactly $t$. Adding an isolated vertex adds a $1\times 1$
identity block to $A+I$, so it increases the complementary $2$-rank by exactly
one and does not create any clique. Hence the resulting graph is still
$K_{m+1}$-free, has complementary $2$-rank exactly $t$, and has at least
\[
(10^h)^{\beta_m-o(1)}
\geq
(t/10)^{\beta_m-o(1)}
=
t^{\beta_m-o(1)}
\]
vertices.

Finally,
\[
\lim_{m\to\infty}
\frac{32}{\left(32+144\cdot 2^{m+1}\right)^{1/(m+1)}}
=
16.
\]
Therefore
\[
\lim_{m\to\infty}\beta_m
=
\frac{\log 16}{\log 10}.
\]
Thus, for every $\varepsilon>0$, one can choose $m$ large enough so that
\[
\beta_m>
\frac{\log 16}{\log 10}-\varepsilon,
\]
which proves the final claim.
\end{proof}

\end{document}
